\documentclass[twocolumn]{notes}
\usepackage{notes}

\begin{document}
\title{Overall plan \& progress}
\date{\today} % \formatdate{10}{11}{2022}
\maketitle
\tableofcontents
\newpage

\section{Big picture}
For ES method in (WTMetaD, OPES, \ldots)
\begin{enumerate}
    \item Data rich CVs
    \begin{enumeratec}
        \item Simulation
        \item Reweight data
        \item Measure sampling progress
        \item Plot sampling progress (different VIZs?)
        \item Evaluate enhancement metric
    \end{enumeratec}
    \item Data poor CVs
    \begin{itemizec}
        \item Chicken and egg: how to enhance sampling using poor quality CVs
        \item How to reweight old data sampled from different bias potential and mix together with new data properly.
        \item Need a way to evaluate the current quality of CVs realative to the ``known'' data rich CVs.
    \end{itemizec}
\end{enumerate}

\subsection{Calculation checklist}
In order to do the above, need to work out:
\begin{enumeratec}
    \item KDE bandwidth issues (accuracy vs. computability) plus fact that CV distributions are NOT Gaussian.  See appendix.
    \item Can VES solve bandwidth problem entirely?
    \item Methods of evaluating
    \begin{itemizec}
        \item ``Quality'' of FES
        \item Enhancement metric
    \end{itemizec}
\end{enumeratec}

\subsection{Immediate priority}
\begin{enumerate}
    \item \textbf{Reweighting, ITRE}.  The weight factors for each frame are split
    \[
        W_t(s_t) \equiv A_t B_t
    \]
    Where the $A_t$ part is frozen (depends on $V_t(s_t)$) and the $B_t$ term contains the (to solve for) bias factor.  Initially $B_t = 1$.
    \item \textbf{KDE compression}.  Look carefully at the bias potential expanded over Gaussians
    \[
        V_n(s) = \sum_k^n h_h\mye^{-\half\norm{s-s_k}^2}
    \]
    When this set of Gaussians is compressed (for the purpose of biased dynamics), how can you use this compressed set of Gaussians to stand in for the KDE representation of the PDF:
    \[
        p(s,n) = \frac{\sum_k^n W_k G(s-s_k)}{\sum_k^n W_k}
    \]
    This will involve pulling parts of the $h_k$ and combining with the weights.  We will need a way to keep count of how many Gaussians have been merged together during compression (reflected in $h_k$) so that the $B_k$ terms can be initialized properly.
    \item \textbf{Validating reweighting algorithms}.  Use measure of quality of
    \begin{enumeratec}
        \item FES
        \item PDF
    \end{enumeratec}
    \item \textbf{Compare enhancement factors}.  Experiment with a number of ways of measuring ``enhancement''.  Is OPES faster than WTMetaD?  Can we get away with our wide bandwidths in both $V_t(s)$ and $p(s, t)$ and get enhancement boost?

    \emph{i.e.} do the forces which are going to be applied crudely in $s$-space (compared to the sharpeness of the peaks in the true $p(s)$ histogram) succeed in driving the system out of the narrow $F(s)$ wells?  This would be a huge win if they do.

    Otherwise, we're stuck with seeking adaptive methods to work with very narrow Gaussians in certain parts of the space.  See Appendix on Silverman's.
\end{enumerate}

\section{ES metrics}
\subsection{Master list}
\begin{enumeratec}
    \item Rate constant-based.  If $\tau = k^{-1}$, Hummer suggests a formula like
    \[
        \rr{Boost} = \sum_i Q_i\frac{\tau_{+,i} + \tau_{-,i}}{\tau_{+,0} + \tau_{-,0}}
    \]
    \item Conformational entropy-based.  Plot $S_\rr{conf}$ over $n$.  Look at the trends.  Compare initial slope.
    \item FES-based.  Plot $\norm{F_\infty-F_n}^2$ vs. $n$ and do the same as above.
    \item PDF-based.  Do the above, except using:
    \[
        p(x, n) = \textrm{estimate of unbiased PDF at step }n
    \]
    \item MFPT-based.  Something like
    \[
        \rr{Boost}_n = \frac{\rr{MFPT}_0}{\rr{MFPT}_n}
    \]
    Recognize that MFPT is directly related to rate constants.
    \item ACF-based.  Generally,
    \[
        C_t = \frac{\langle(f(x_t)-\bar{f})f(x_0)-\bar{f})}{\myVar{f(x)}}
    \]
    The most sensible observables for which to inspect autocorrelation are the eigenfunctions which are zero-mean and unit variance by construction.  Use the known data-rich functions.  Ah!  They are not unit variance over limited datasets!
    \[
        C_t \Leftarrow \{\psi_1(s_k),\,\psi_1(s_{k+t})\}_{k=1}^n
    \]
    \item ``Best practise'' = measure distances for FES/PDF and do a number of ES simulations to establish the central tendency and variance of convergence rates.
    \item ``Best practise'' = for FES, bootstrap error bars to observe when convergence has occured to within uncertainty.
    \item Exploration-extent based.
    \begin{itemizec}
        \item fraction of alanine dipeptide states visited (\emph{e.g.} 2/6)
        \item fraction of $\Omega_s$ explored.  Area or fraction of non-zero histogram bins.  This is going to depend on bin-size, so is a little rough.
    \end{itemizec}
\end{enumeratec}
A big idea here is to try a bunch/all of these methods and see if you identify any trends or if there is something else instructive.

\subsection{Conformational entropy}
This is
\[
    S = -k_\mathrm{B}\int\myd{x}\rho(x)\log\rho(x) =
    +k_\mathrm{B}\int\myd{x}\rho(x)\log\frac{1}{\rho(x)}
\]
\begin{figure}[h!]\centering    %%  caption below figure
  \includegraphics[width=0.4\textwidth]{plogp.png}
  % \caption{General caption}
  % \label{fig:label}
\end{figure}
The plot shows how the integrand depends on the density $p(x)$.  Computationally, for a 2D density,
\[
    S = k_\mathrm{B}\Delta x_1\Delta x_2 \sum_{i,j}^N
    \rho(x_{ij})\log\frac{1}{\rho(x_{ij})}
\]
This entropy will have a maximum when
\[
    \rho(x) = Z^{-1}\mye^{-\beta E(x)}
\]
This can be shown by optimizing
\[
    G = E - TS \quad\Rightarrow\quad \myfunctional{G}{\rho} = 0
\]
subject to the constraint that $\rho$ is normalized.

Another way to approach conformational entropy is to follow Karplus and break the conformational space into $n$ distinct metastable states.  Within each state $s_i$ there will be some internal conformational entropy, $S_i^b$.  There is flexibility in how this is computed, whether as a sub-histogram (re-normalized for each free energy basin) or by carrying out Normal Mode Analysis for an energy basin.

In ths way, the fraction of being in any discrete state is given by $p_i$ and the conformational entropy of the system is
\[
    S = \sum_i^n p_i\,S_i^b - k_\mathrm{B}\sum_i^n p_i\log p_i
\]
\subsection{2D surfaces: PDFs and FESs}
For convenience, define any surface
\[
    u(x, y)
\]
over $N\times N$ grids so that integrating over the surface is
\[
    I = \int_{x_0}^{x_N}\myd{x}\int_{y_0}^{y_N}\myd{y}u(x, y)
    \sim \Delta x\Delta y\sum_{i,j}^Nu(x_i,y_j)
\]
and the domain has volume
\[
    |\Omega| = (x_N-x_0)(y_N-y_0) = N^2\Delta x\Delta y
\]
The standard gotcha:  A grid containing $N$ points is defined by $(N+1)$ bin edges.  Taking the Riemann sum defined by the function/surface at the the left (lower) bin edge give the following definition of the mean.
\[
    \bar{u} = \frac{1}{|\Omega|}\int\dd x\int\myd{y}u(x,y)
    \sim \frac{1}{N^2}\sum_{i,j}^N u(x_i,y_j)
\]

\subsection{Distance of a histogram from a reference}
Recommended to use Jensen-Shannon divergence.  This is due to computational/mathematical considerations about ratios of histograms where the support of the numerator extends to places where the denominator is 0, which blows up the sum.
\[
    D_\rr{JS}(P\parallel Q) = \half\Big[ D_\rr{KL}(P\parallel M) + D_\rr{KL}(Q\parallel M)\Big]
\]
where $M$ is the intermediate distribution
\[
    M = \half [P + Q]
\]
and $D_\rr{KL}$ is
\begin{align*}
    D_\rr{KL}(P\parallel Q) &= \int\dd x\int\myd{y}p(x,y)\log\frac{p(x,y)}{q(x,y)} \\
    &\sim \Delta x\Delta y\sum_{i,j}^N p(x_i,y_j)\log\frac{p(x_i,y_j)}{q(x_i,y_j)}
\end{align*}

\subsection{Distance of a FES from a reference}
Either MSE and RMSD are very similar.  For a 2D FES sampled up to step $n$ (relative to the data rich FES, $F_\infty(x)$):
\begin{align*}
    \rr{MSE}(F_\infty, F_n) &= \frac{1}{|\Omega|}\int\dd x\int\myd{y}
    \Big[F_\infty(x,y) - F_n(x,y) \Big]^2 \\
    &\sim \frac{1}{N^2} \sum_{i,j}^N
    \Big[F_\infty(x_i,y_j) - F_n(x_i,y_j) \Big]^2
\end{align*}
while the RMSD is (to a factor involving the number of bins and the volume of $x$-space)
\[
    \rr{RMSD}(F_\infty, F_n) = \sqrt{\rr{MSE}(F_\infty, F_n)}
\]

\section{Reweighting, ITRE}
The correct weights for reweighting ES data from a time-dependent bias $V_t(s)$ is found by writing out the desired ensemble $p(s)$ and the instantaneous biased $p_t(s)$, time $t$
\[
    W_t(s_t) = \frac{Z_t}{Z}\mye^{+\beta V_t(s_t)}
    = \mye^{+\beta[V_t(s_t)-c_t]}
\]
This defines the so-called bias shift $c_t$.  The key observation is that the first part of the weight is simple and permanently fixed
\[
    A_t := \mye^{+\beta V_t(s_t)}
\]
The difficulty is in solving for
\[
    B_t := \mye^{-\beta c_t} = \frac{Z_t}{Z} = \langle \mye^{-\beta V_t} \rangle
\]
The average above is over the unbiased ensemble.  The basic idea of ITRE is to iteratively solve for them using the definition for the KDE estimate for $p(s)$ at time $t$
\[
p(s,t) = \sum_{t'}^t \widetilde{W}_{t'} G(s-s_{t'}),\quad
\widetilde{W}_{t'} = \frac{W_{t'}}{\sum_{t''}^{t}W_{t''}}
\]
Because in metadynamics $V_t$ is changing so often, there will be many $B_t$ terms to estimate
\[
    B_t = \int\myd{s}\tilde{p}(s)\mye^{-\beta V_t(s)},\quad
    \tilde{p}(s) = \frac{p(s)}{\int\myd{s}p(s)}
\]
The trick is to use the current best estimate for $p(s)$, $p(s,t)$ in this expression and then use theee terms to construct a new set of weights $W_t$.  Iterate until self-consistent.

\subsection{Explore efficient computation of Gaussian integrals}
Combining the two ITRE equations (and assuming that the KDE representation si properly normalized, \emph{e.g.} ignoring the renormalization) gives
\[
    B_T = \frac{\sum_{t=0}^T A_tB_t G_{tT}}{\sum_{t=0}^T A_tB_t}
\]
with
\[
    G_{tT} = \int\myd{s}\mye^{-\beta V_T(s)}G(s-s_t)
\]
So the question is, is there anything here?  Is there a tractable way of efficiently computing integrals of this type?  MANY of them would need computing
\[
    n\times T
\]
where $T$ is the total number of frames used in the KDE representation of $p(s)$ and $n$ is the total number of metadynamics bias potentials.

A possibility is that KDE compression can reduce this number.

A possibility is that the $V_t(s)$ functions being Gaussians themselves might lead to something clever.  ChagGPT says no.

ChatGPT recommends, for integrals of the form
\[
    \int\myd{x}\mye^{-f(x)}G(x-x_0)
\]
\begin{enumeratec}
    \item Laplace/saddle point
    \begin{itemizec}
        \item Laplace approximation
        \item Gaussian integral with nonquadratic potential
        \item MacKay-Laplace approximation
    \end{itemizec}
    \item Delta method (see notes/plan/gaussian-.pdf)
\end{enumeratec}

\subsection{Explore nature of $B_t$ integral}
Main practical issue in
\[
    c_t = \kT\log\myE{\mye^{-\beta V_t(s)}}
\]
\begin{enumeratec}
    \item variance/numerical stability of expectation/avg.
    \item how to re-use computational work across set of $V_t$
\end{enumeratec}
ChatGPT suggests
\begin{enumeratec}
    \item stable evaluation, log-sum-exp trick.  log-mean-exp.  ``Numerically stable free energy''  ``Cumulant generating log-moment''
    \item Free energy
    \begin{itemizec}
        \item Zwanzig exponential averaging (free energy perturbation)
        \item Jarzynski equality has same structure
    \end{itemizec}
    \item Monte carlo
    \begin{itemizec}
        \item quasi-Monte carlo (QMC)
        \item stratified sampling / mixture importance sampling
        \item control variates
    \end{itemizec}
    \item Density estimation and integration.  ``Importance sampling using emiracle distribution.''
\end{enumeratec}
ChatGPT: ``Best first option'' = data representation (weighted MC average).  This would be to compute and store
\[
    C_{tt'} = \mye^{-\beta V_t(s_{t'})}
\]
For each new frame $s_T$, must compute full set of potentials
\[
    \{\mye^{-\beta V_t(s_T)}\}_{t=0}^T \quad \{V_t(s_T)\}_{t=0}^T
\]

\subsection{Representing PDFs: exact $p(s)$, $B_t$}
This is needed to benchmark different approaches for computing the weights.

The representation of $p(s)$ is of central importance.  There are some options
\begin{enumeratec}
    \item histogram (should be equivalent to MC, so can skip)
    \item KDE (need to explore bandwidths)
    \item GMM (need to explore scikit-learn Bayesian Gaussian Mixture)
    \item GMM / MC (sample the GMM)
    \item Gaussian processes
\end{enumeratec}

ChatGPT says
\begin{enumeratec}
    \item weighted data representation = ``weighted empiracle measure''
    \[
        \langle A\rangle = \sum_k \tilde{W}_k A(x_k)
    \]
    \item KDE: recommend if you intend to do grid-based quadrature
    \begin{itemizec}
        \item sensitive to bandwidth
        \item robust non-parametric method
        \item expensive scaling with data size
    \end{itemizec}
    \item GMM: compact, easy to sample
    \begin{itemizec}
        \item must choose $K$
        \item may struggle with sharp peaks if $K$ too small
        \item excellent if $A(x)$ gives analytical integrals
    \end{itemizec}
    \item Normalizing flow
    \begin{itemizec}
        \item fast sampling once trained
        \item need lots of data for training
        \item well-separated modes are a challenge
        \item overkill for 2D
    \end{itemizec}
\end{enumeratec}


\subsection{Validating FESs}
\begin{enumeratec}
    \item barrier heights
    \item relative well heights
\end{enumeratec}

\subsection{Integrals: grids vs. MC}
\textbf{Adaptive quadrature}: in the case that you're doing any integrals over grids, look into existing python packages for this.  Look at cost-savings over naive grids.  ``Deterministic quadrature on adaptive meshes.''

Is MC the way to go?  Not likely.  Look into QMC (Sobol-Halton).  The structure of $B_t$ makes for a challenging avg. to sample/converge.

\subsection{TODO}
PDFs, exact $p(s)$
\begin{itemizec}
    \item scikit learn Bayesian Gaussian Mixture
    \item investigate Gaussian process package
    \item investigate metrics to compare PDF representations
    \item other?  Look at recent research
    \item explore mathematical connection between simple KDE and ``importance sampling using the empiracle distribution''.  They must be equivalent in the limit of delta function.
    \item QMC (Sobol-Halton)
\end{itemizec}
Computing ITRE terms
\begin{itemizec}
    \item Work out reconstructing $V_t$ from WTMetaD simulation data
    \item Work out exact $B_t$ terms using these $V_t$
    \item Look at typical size distributions of $A_t$, $B_t$, $G_{tt'}$ terms
    \item Suppose there is a way to compute $G_{tt'}$ terms.  Implement the algorithm as it stands and compare to ``exact'' $B_t$ terms.
    \item KDE compression in reducing the number of $G_{tt'}$ terms (or computational effort in computing individual ones)
    \item ChatGPT:  feed it the 2020, Giberti ITRE paper and think up some good questions.  E.g. sort out the renormaization of $p(s)$
\end{itemizec}

\section{Reweighting, WHAM}
I am doing enhanced sampling simulations of molecular systems.  There are a number of scenarios:
\begin{enumerate}
    \item Simgle unbiased ensemble
    \[
        \langle O\rangle = \int\myd{s}p(s)\,O(s)
        =\lim_{N\rightarrow\infty} \frac{1}{N}\sum_{k=1}^N O(s_k), ~s_k\sim p
    \]
    Unbiased distribution
    \[
    p(s) = \frac{\mye^{-\beta F(s)}}{Z}
    \]
    Biased distribution with time-dependent bias potential:
    \[
        p_t(s) = \frac{\mye^{-\beta [F(s)+ w_t(s)}}{Z_t}
    \]
    Reweighting factor for Monte Carlo sampling using biased simulation data:
    \[
        W_t(s_t) = \mye^{\beta[w_t(s_t)-c_t]} ,~~~\mye^{-\beta c_t} = \frac{Z_t}{Z}
    \]
    Estimating factors involving $c_t$ using biased simulation data
    \[
        \mye^{-\beta c_t} \approx \frac{\sum_{t'=1}^T W_{t'}(s_{t'})\mye^{-\beta w_t(s_{t'})}}{\sum_{t'=1}^T W_{t'}(s_{t'})}
    \]
    where the weights themselves depend on $c_t$.
\end{enumerate}

\subsection{Connection Between Reweighting and WHAM}

The key point is that there are two distinct kinds of weights:

\begin{enumerate}
    \item the \emph{single-ensemble reweighting factor}, and
    \item the \emph{multi-ensemble WHAM combination weights}.
\end{enumerate}

These serve different purposes.

\subsection{Single-Ensemble Reweighting}

Suppose the unbiased distribution is

\[
p(s)=\frac{e^{-\beta F(s)}}{Z}.
\]

A biased ensemble with time-dependent bias potential \(w_t(s)\) samples

\[
p_t(s)=\frac{e^{-\beta[F(s)+w_t(s)]}}{Z_t}.
\]

Define the reweighting factor

\[
W_t(s)=e^{\beta[w_t(s)-c_t]},
\]

with

\[
e^{-\beta c_t}=\frac{Z_t}{Z}.
\]

Then

\[
p(s)=p_t(s)e^{\beta[w_t(s)-c_t]}.
\]

Thus, the factor

\[
e^{\beta[w_t(s)-c_t]}
\]

removes the umbrella bias and converts samples from the biased ensemble into unbiased samples.

This is exactly the same structure as Roux Eq.~(7),

\[
[(p(\xi))]^{\mathrm{unbiased}}_{(i)}
=
e^{+\beta w_i(\xi)}
(p(\xi))_{(i)}
e^{-\beta F_i}.
\]

The correspondence is

\[
F_i \leftrightarrow c_i.
\]

Indeed, Roux defines

\[
e^{-\beta F_i}
=
\left\langle e^{-\beta w_i}\right\rangle
=
\frac{Z_i}{Z},
\]

which is identical to the normalization ratio above.

At this stage, we are still dealing with a \emph{single} ensemble.

\subsection{WHAM Introduces a Second Weighting Layer}

Suppose we now have many umbrella windows \(i=1,\dots,K\).

Each window produces an unbiased estimator

\[
\hat p_i(s)
=
p_i(s)e^{\beta[w_i(s)-F_i]}.
\]

If sampling were infinite, all \(\hat p_i(s)\) would agree exactly.

However, with finite sampling:

\begin{itemize}
    \item some windows are noisy,
    \item some regions are poorly sampled,
    \item overlap between windows may be uneven.
\end{itemize}

WHAM asks:

\[
\textit{How should all these estimators be combined optimally?}
\]

\subsection{The WHAM Equation}

Roux Eq.~(6) gives

\[
p(s)
=
\frac{
\displaystyle
\sum_i n_i p_i(s)
}{
\displaystyle
\sum_j n_j e^{-\beta[w_j(s)-F_j]}
}.
\]

This can be rewritten as

\[
p(s)
=
\sum_i a_i(s)\,\hat p_i(s),
\]

where

\[
a_i(s)
=
\frac{
n_i e^{-\beta[w_i(s)-F_i]}
}{
\displaystyle
\sum_j n_j e^{-\beta[w_j(s)-F_j]}
}.
\]

The factors \(a_i(s)\) are the WHAM combination weights.

Thus there are \emph{two} layers:

\begin{enumerate}
    \item Bias removal:
    \[
    e^{\beta[w_i(s)-F_i]},
    \]
    which converts a biased distribution into an unbiased one.

    \item Statistical combination:
    \[
    a_i(s),
    \]
    which combines all windows into a minimum-variance estimate.
\end{enumerate}

\subsection{Interpretation of the WHAM Weights}

Suppose \(s\) lies near the center of umbrella \(i\). Then

\[
w_i(s)\approx 0,
\]

so umbrella \(i\) samples that region efficiently.

For another umbrella \(j\) far away,

\[
w_j(s)\gg 1,
\]

and that window contributes very little statistical information there.

Consequently,

\[
a_i(s)
\propto
n_i e^{-\beta[w_i(s)-F_i]}
\]

gives large weight to windows that sample \(s\) well.

Thus:

\begin{itemize}
    \item the reweighting factor removes the bias,
    \item the WHAM factor determines statistical confidence.
\end{itemize}

These are conceptually distinct.

\subsection{Relation to Multi-Ensemble Monte Carlo}

WHAM is fundamentally a multi-ensemble importance sampling method.

Each umbrella window samples

\[
\pi_i(x)
\propto
e^{-\beta[U(x)+w_i(x)]}.
\]

The unknown constants \(F_i\) are relative free energies,

\[
F_i
=
-k_B T \ln \frac{Z_i}{Z}.
\]

WHAM determines the \(F_i\) self-consistently so that all windows agree on the same underlying unbiased distribution.

This is closely related to:

\begin{itemize}
    \item Bennett acceptance ratio (BAR),
    \item Ferrenberg--Swendsen reweighting,
    \item MBAR (multistate Bennett acceptance ratio).
\end{itemize}

In fact:

\[
\begin{aligned}
\text{BAR} &\rightarrow \text{two ensembles}, \\
\text{WHAM} &\rightarrow \text{histogram formulation}, \\
\text{MBAR} &\rightarrow \text{continuous/unbinned generalization}.
\end{aligned}
\]

\subsection{Answer to the Original Question}

Yes.

Roux Eq.~(7) gives the unbiased estimator from a single window.

WHAM then introduces an additional weighting factor to combine all windows optimally.

Schematically:

\[
\text{biased histogram}
\quad
\overset{\text{Eq.~7}}{\longrightarrow}
\quad
\text{unbiased local estimate}
\quad
\overset{\text{WHAM}}{\longrightarrow}
\quad
\text{global optimal estimate}.
\]

Thus the WHAM weights are an additional layer on top of the single-window reweighting factors.

\section{Assessing quality of CVs}
Anything coming off an SRV will have
\[
    \myE{\psi_i} = 0 \quad \myE{\psi^2_i} = 1 \quad \psi_i=\sum_j a_{ij}\zeta_j
\]
Work out what this means for when you evaluate the quality of $\psi_i$ when evaluated using ACF.  Will have to validate using full unbiased dataset.

Contemplate the meaning of
\[
    \tau_i \Leftarrow \myE{\psi_i(t+\tau)\psi_i(t)}
\]
and take care about the denominator in the usual ACF.

\subsection{Advances in KDEs and related methods}
From ChatGPT
\begin{enumeratec}
    \item Adaptive bandwidth.  Abramson (1982)
    \[
        \sigma_i \propto\frac{1}{\sqrt{p(x_i)}}
    \]
    Regions that are highly sampled get fine structure.  Silverman (1986) gives low bandwidth will increasing effective sample size.  Generally, bandwidth starts large and gets small.
    \item Anisotropy.  Terrell and Scott (1992).  Boter et al (2010).  Chacon and Duong (2018).
    \item kNN-based.  Choose $\sigma_i$ based on nearest neighbour.
    \item GMM.  NOT KDE.  Karplus, 2009

    See https://pubmed.ncbi.nlm.nih.gov/19284746
    \item Gaussian processes for FES.  Mones, 2016.  ``Exploration, Sampling, And Reconstruction of Free Energy Surfaces with Gaussian Process Regression''
    \item Normalizing flows.  Noe, 2019.  Not a priority for my 2D FES: stuggles with well-separated peaks, overkill for 2D, we lack a priori knowledge of $F(s)$ (Noe leverages Boltzmann factor).
\end{enumeratec}

\subsection{TODO}
\begin{itemize}
    \item review Noe 2019 and look into whether it has a place with a low-dimensional (2D) FES compared to the PES he uses.  I think he leverages
    \[
        p(x) \sim \mye^{-\beta V(x)}
    \]
    to learn the transformation to a $z$ variable that is $3N$ (high) dimensional but with simple structure/shape.  We are trying to learn an FES, so we can similarly leverage knowlege of $F(s)$.
\end{itemize}

\section{Search}
\begin{itemizec}
    \item "2D kernel density estimation" astronomy large dataset
    \item "adaptive kernel density estimation" sharp peaks 2D
    \item "bivariate density estimation" cross-validation bandwidth
    \item "least squares cross validation" kernel density estimation 2D
    \item "spatial hotspot" kernel density estimation bandwidth
    \item "home range estimation" kernel density estimation LSCV
    \item "bivariate KDE" multimodal density estimation
    \item "density estimation" "sharp peaks" "2D"
\end{itemizec}


\section{WT-Metadynamics key results and the PLUMED bias shift}

The reasoning here begins by proposing to sample a free energy surface that has desireable qualities:
\[
    p_\gamma(s) \propto [p(s)]^{1/\gamma} ,\quad \gamma = \frac{T+\Delta T}{T}
\]
This is to be accomplished by biasing the simulation with a some potential $V_t(s)$.  So the question is then, what bias will give us our target distribution?
\[
    \frac{\mye^{-\beta[F(s)+V_t(s)]}}{Z_t} = \frac{\mye^{-\beta F(s)/\gamma}}{Z_\gamma}
\]
or, introducing the free energies associated with each ensemble
\[
    F(s)+V_t(s)-F_t = \gamma^{-1}F(s) - F_\gamma
\]
This gives
\[
    F(s) = [\gamma^{-1}-1]^{-1}\Big(V_t(s) - [F_t-F_\gamma] \Big)
\]
Now clean up
\[
    [\gamma^{-1}-1]^{-1} = \frac{1}{\gamma^{-1}-1} = \frac{\gamma}{1-\gamma}
\]
\textbf{Question}: what is the physical meaning of $F_t-F_\gamma$?  Does this become zero as the ensemble converges asymptotically toward the well-tempered ensemble?  Clearly, if they are the same ensemble, the MUST have the same free energy?

Ok, we now have
\begin{align*}
    \beta F(s) &= \beta\frac{\gamma}{1-\gamma}V_t(s) - \beta\frac{\gamma}{1-\gamma}[F_t-F_\gamma] \\
    &= \gamma\frac{\beta}{1-\gamma}V_t(s) - \gamma\frac{\beta}{1-\gamma}\Delta F_{\gamma t} \\
    &= \gamma\Delta\beta V_t(s) - \gamma\Delta\beta \Delta F_{\gamma t}
\end{align*}
defining
\[
    \frac{\beta}{1-\gamma} = \Delta\beta
\]
This expression for $\beta F(s)$ above comes indirectly.  We asked, what bias gives the WT ensemble?  The answer involves the canonical ensemble and $F(s)$.  Rearranging the expression for $V_t(s)$ for $F(s)$ allows substitution back in the definition of the bias shift and gives the PLUMED formula for it.  We're looking at
\[
    \mye^{-\beta c_t} = \frac{Z_t}{Z} = \frac{\int\myd{s}\mye^{-\beta[F(s)+V_t(s)]}}{\int\myd{s}\mye^{-\beta F(s)}}
\]
Now, conceptually what we're looking at here is saying IF $V_t(s)$ gives the well-tempered ensemble, it can be written in terms of $F(s)$.  We do the converse.  Write $F(s)$ in terms of the proper, well-tempered bias $V_t(s) = V_\gamma(s)$:
\begin{align*}
    \beta[F(s)+V_t(s)] &= \gamma\frac{\beta}{1-\gamma}V_t(s) - \gamma\Delta\beta\Delta F_{\gamma t} + \beta V_t(s) \\
    &= \beta\left[\frac{\gamma}{1-\gamma}+1\right]V_t(s) - \gamma\Delta\beta\Delta F_{\gamma t} \\
    &= \beta \frac{1}{1-\gamma}V_t(s) - \gamma\Delta\beta\Delta F_{\gamma t} \\
    &= \Delta\beta V_t(s) - \gamma\Delta\beta\Delta F_{\gamma t}
\end{align*}
The time-dependent $\gamma\Delta\beta\Delta F_{\gamma t}$ terms cancel in the partition function ratio, and the PLUMED expression is
\[
    \mye^{-\beta c_t} = \frac{Z_t}{Z} = \frac{\int\myd{s}\mye^{-\Delta\beta V_t(s)}}{\int\myd{s}\mye^{-\gamma\Delta\beta V_t(s)}}
\]

\subsection{Connection to Tiwary, 2015, eq. 19}
Tiwary's expression is
\begin{align*}
    F(s) &= -\left(\frac{\gamma}{\gamma-1}\right)V_t(s) + \frac{1}{\beta}
    \log\int\myd{s}\mye^{\gamma\Delta\beta V_t(s)} \\
    \beta F(s) &= +\gamma\Delta\beta V_t(s) + \frac{1}{\beta}
    \log\int\myd{s}\mye^{\gamma\Delta\beta V_t(s)}
\end{align*}


\appendix
\section{Dealing with a breakdown of Silverman's}

\subsection{Adaptive / local bandwidth KDE}
Instead of a global $\sigma$, use position-dependent bandwidths.  Key ideas:
\begin{itemizec}
    \item Narrow kernels in dense regions (sharp peaks)
    \item Wide kernels in sparse regions (barriers)
    \item Examples in literature:
    \item Adaptive KDE (balloon / sample-point estimators)
    \item k-nearest-neighbor bandwidth selection
    \item Diffusion KDE
\end{itemizec}
Follow-ups to OPES (e.g. OPES-explore / OPES-metadynamics variants).

Allows:
\begin{itemizec}
    \item per-kernel covariance
    \item metric-based scaling
\end{itemizec}
BUT, method is still not fully “standardized” in mainstream OPES workflows.

\subsection{Mixture models instead of KDE}
Replace KDE with something like:
\begin{itemizec}
    \item Gaussian Mixture Models (GMM)
    \item Variational density estimators
\end{itemizec}
This idea appears earlier in:
\begin{itemizec}
    \item Gaussian-mixture umbrella sampling (Maragakis et al.)
    \item Later ML-based approaches
\end{itemizec}
Benefits
\begin{itemizec}
    \item Explicitly models multimodality
    \item Each mode gets its own covariance
\end{itemizec}
Other approaches
\begin{itemizec}
    \item Variational Enhanced Sampling (VES)
    \item Neural density estimators (normalizing flows)
\end{itemizec}
Tradeoff:
\begin{itemizec}
    \item Requires optimization
    \item Less “plug-and-play” than OPES
\end{itemizec}

\subsection{Biasing without explicit density estimation}
Methods:
\begin{itemizec}
    \item Variational Enhanced Sampling (VES)
    \item Neural network biases
    \item Reinforcement learning biasing
\end{itemizec}
This handles the bandwidth problem and can represent sharp features naturally.
It is more complex, however, and requires basis choice / neural architecture.

\subsection{Density-dependent bandwidth (Abramson-type)}
Shrink bandwidth where density is high:
\[
    \sigma_k \propto P(s_k)^{-\alpha} \quad \alpha \approx \half
\]

\subsection{k-nearest-neighbor (kNN) bandwidth}
Idea = use bandwidth based on nearest neighbour bandwidth.  Something to look into.  Less intersting than GMM. Appears in:
\begin{itemizec}
    \item Bayesian reconstruction methods
    \item Some Gaussian process / regression-based FES methods
    \item Kernel compression literature
\end{itemizec}

\subsection{Adaptive Gaussian MetaD}
Replace OPES altogether.  2012, Branduardi.  MetaD Gaussians

\subsection{Machine learning density models}
Increasingly used:
\begin{itemizec}
    \item Normalizing flows
    \item Neural density estimators
\end{itemizec}

\section{What is popular now?}
\begin{enumeratec}
    \item Bias-potential / CV-based methods (still dominant)
    \begin{itemizec}
        \item Metadynamics (and variants)
        \item OPES
        \item Variational Enhanced Sampling (VES)
    \end{itemizec}
    \item  Replica / tempering methods
    \begin{itemizec}
        \item Replica Exchange MD (REMD, REST2, etc.)
    \end{itemizec}
    \item Path / ensemble methods (unbiased or weakly biased)
    \begin{itemizec}
        \item Weighted Ensemble (WE)
        \item Transition Path Sampling (TPS)
        \item Adaptive sampling (MSM-driven)
    \end{itemizec}
    \item Machine-learning–driven methods (fastest growing)
    \begin{itemizec}
        \item Learned CVs
        \item Neural bias potentials
        \item Reinforcement learning / flows
    \end{itemizec}
\end{enumeratec}

\subsection{Modern metadynamics}
\begin{enumeratec}
    \item Transition-tempered / adaptive MetaD
    \item Parallel bias / multiple walkers
    \item Better reweighting.  2023, Stevensson \& Eden
    \item Multiwalker VES.
\end{enumeratec}

\end{document}



%% \href{https://www.ferglab.com/}{\color{darkred}Website}

% \usepackage{minted}
% \begin{minted}[mathescape, linenos]{python}
%     # Note: $\pi=\lim_{n\to\infty}\frac{P_n}{d}$
%     title = "Hello World"
% \end{minted}

%%=======================================================================
%%     PREFERRED FIGURE/TABLE
%%=======================================================================
%%                                 %%-------------------------------------%%
%% \begin{figure}[h!]              %%  caption right of figure
%%   \begin{minipage}[c]{0.67\textwidth}
%%     \includegraphics[width=\textwidth]{imgfile.png}
%%   \end{minipage}\hfill
%%   \begin{minipage}[c]{0.3\textwidth}
%%     \caption{Caption on the right hand side}
%%     \label{fig:}
%%   \end{minipage}
%% \end{figure}
%%                                 %%-------------------------------------%%
%% \begin{figure}[h!]\centering    %%  caption below figure
%%   \includegraphics[width=0.5\textwidth]{imgfile.png}
%%   \caption{General caption}
%%   \label{fig:label}
%% \end{figure}
%%                                 %%-------------------------------------%%
%% \begin{figure}[!h]\centering    %%  array of figures with subcaptions
%%   \subfloat[]{\includegraphics[width=0.5\textwidth]{}}
%%   \subfloat[]{\includegraphics[width=0.5\textwidth]{}} \\
%%   \subfloat[]{\includegraphics[width=0.5\textwidth]{}}
%%   \subfloat[]{\includegraphics[width=0.5\textwidth]{}} \\
%%   \subfloat[]{\includegraphics[width=0.5\textwidth]{}}
%%   \subfloat[]{\includegraphics[width=0.5\textwidth]{}}
%%   \caption{General caption}
%%   \label{fig:label}
%% \end{figure}
%%                                        %%------------------------------%%
%% \begin{table}[!h]\centering\ra{1.3}    %%  simple table
%%   \begin{tabular}{@{}rrrr@{}}\toprule
%%     & $t=0$ & $t=1$ & $t=2$ \\ \midrule
%%     $c$ & 1 & 2 & 3 \\
%%     $c$ & 1 & 2 & 3 \\
%%     \bottomrule
%%   \end{tabular}
%%   \caption{Caption}
%% \end{table}
%%                                        %%------------------------------%%
%% \begin{table}[!h]\centering\ra{1.3}    %%  complicated table
%%   \begin{tabular}{@{}rrrrcrrr@{}}\toprule
%%     & \multicolumn{3}{c}{$w = 8$} & \phantom{abc}&
%%     \multicolumn{3}{c}{$w = 16$}\\
%%     \cmidrule{2-4} \cmidrule{6-8} \cmidrule{10-12}
%%     & $t=0$ & $t=1$ & $t=2$ && $t=0$ & $t=1$ & $t=2$ \\ \midrule
%%     \rowcolor[gray]{.95} $c$ & 1 & 2 & 3 && 4 & 5 & 6 \\
%%     \rowcolor[gray]{1.0} $c$ & 1 & 2 & 3 && 4 & 5 & 6 \\
%%     \bottomrule
%%   \end{tabular}
%%   \caption{Caption}
%% \end{table}

%%=======================================================================
%%     OLDER FIGURE/TABLE
%%=======================================================================
%% \begin{center}
%%   \includegraphics[width=0.75\linewidth]{}
%%   \captionof{figure}{}
%% \end{center}
%%
%% \begin{center}
%%   \begin{tabular}{rl}
%%     \includegraphics[width=0.5\linewidth]{} &
%%     \includegraphics[width=0.5\linewidth]{} \\
%%     \includegraphics[width=0.5\linewidth]{} &
%%     \includegraphics[width=0.5\linewidth]{} \\
%%   \end{tabular}
%%   \captionof{figure}{}
%% \end{center}
%%
%% \begin{center}
%%   \begin{tabular}{lll}
%%     \hline
%%     A & B & C \\
%%     \hline
%%     \hline
%%     1 & 2 & 3 \\
%%     \hline
%%   \end{tabular}
%%   \captionof{table}{test}
%% \end{center}
